\documentclass{situationist-preprint}

\title{Uncertainty and Revolutionary Action}
\msid{POL-23-b}
\msfamily{POL-23}
\msdiscipline{Political Philosophy}
\mstype{Companion manuscript}
\mskeywords{revolution; uncertainty; moral luck; evidence-relative reasons; civil resistance}

\begin{document}
\maketitle

\begin{abstract}
POL-23-a argues that revolution against extreme injustice can be justified, but that its justification depends on whether the resulting order improves on the old one by a margin proportionate to the costs, which cannot be known when the revolution begins. This leaves the agent who must decide whether to take part with a question the principal answers only in outline. We argue that reasonable participation has three conditions: evidence that the order is extremely unjust, a credence in sufficient improvement that rises with the costs the revolution threatens, and a choice of means that limits the harm done if the revolution fails. The third condition gives the agent some control over the moral luck that the principal identifies, and favors methods of civil resistance where they are available.
\end{abstract}

\section{From justification to reasonableness}

Definition~2.3 of POL-23-a distinguishes the justification of a revolution, which depends on the facts including its outcome, from the reasonableness of taking part, which depends on the evidence available to the agent. Theorem~3.2 shows that justification cannot be known at the outset, and Corollary~3.3 that participation may be reasonable although the revolution turns out unjustified. The distinction is that between fact-relative and evidence-relative reasons \citep{Parfit2011}. What it leaves open is when participation is reasonable.

\begin{proposition}\label{p:conditions}
Taking part in a revolution is reasonable for an agent only if (R1) her evidence supports the judgment that the order is extremely unjust in the sense of Definition~2.2 of POL-23-a; (R2) her credence that the revolution will satisfy the condition of Principle~P2 of POL-23-a is high enough, given the costs that it threatens if it fails; and (R3) among the means available to her, she chooses those whose costs are least on the assumption that the revolution fails.
\end{proposition}

\begin{proof}
(R1): By Principle~P1 of POL-23-a, extreme injustice is necessary for justification, and unlike the outcome it is a present fact, assessable on present evidence; an agent without evidence of it has no expected reason that could outweigh the costs. (R2): Participation is reasonable when the expected reasons for it outweigh the expected reasons against (Definition~2.3 of POL-23-a). The reasons against are the costs of revolution, which are borne in any case and are greatest when it fails; so the greater those costs, the higher the probability of success needed for the balance to favor participation. (R3): Of two courses that are equally likely to succeed, the one that does less harm in the event of failure has the greater balance of expected reasons. So a reasonable agent prefers it.
\end{proof}

\section{Means and moral luck}

\citet{Williams1981} argued that some agents' decisions are justified, or not, by outcomes they could not foresee, and POL-23-a treats the outcome-dependence of revolution as an instance. Condition R3 shows that the agent's exposure to such luck is not fixed. She cannot control whether the revolution succeeds, but she can affect how much harm her participation does if it does not, and thereby how much of the justification of her conduct is left to the outcome.

This has a consequence for the choice of methods. Campaigns of nonviolent resistance impose lower costs on others when they fail than armed campaigns do. There is also evidence that, in the twentieth and early twenty-first centuries, they succeeded more often and were more often followed by stable democratic institutions \citep{ChenowethStephan2011}. If that evidence is sound, nonviolent methods improve the agent's position under R2 and R3 together. The conclusion is comparative: where nonviolent methods are unavailable, as under a regime that answers them with massacre, R3 selects among the remaining means, and armed resistance may be the least costly.

\begin{objection}
On this account no one could reasonably have begun a revolution that history now judges justified, since at the outset the evidence for its success was weak.
\end{objection}

\begin{reply}
The objection overlooks that R2 is relative to the costs at stake. Under extreme injustice the cost of not acting is also great, and it enters the balance: when the existing order kills and imprisons on a large scale, a modest probability of success can suffice. Locke's observation that a people will not rise for small mismanagement, but only for a long train of abuses, describes the conditions under which R1 is met and the balance under R2 shifts \citep[\S225]{Locke1690}.
\end{reply}

\section{Conclusion}

That the justification of revolution depends on outcomes that cannot be known in advance does not leave the agent without guidance. She can know whether the injustice is extreme, she can proportion her confidence to the stakes, and she can choose means that limit the harm done if she is wrong. This develops the conclusion of POL-23: revolution is justified under conditions of extreme injustice, though its justification depends on outcomes that cannot be known in advance.

\begin{thebibliography}{9}\small
\bibitem[Chenoweth and Stephan(2011)]{ChenowethStephan2011} Chenoweth, E. and Stephan, M. J. (2011). \emph{Why Civil Resistance Works: The Strategic Logic of Nonviolent Conflict}. Columbia University Press.
\bibitem[Locke(1690)]{Locke1690} Locke, J. (1690). \emph{Two Treatises of Government}. Awnsham Churchill.
\bibitem[Parfit(2011)]{Parfit2011} Parfit, D. (2011). \emph{On What Matters}, vol.~1. Oxford University Press.
\bibitem[Williams(1981)]{Williams1981} Williams, B. (1981). \emph{Moral Luck: Philosophical Papers 1973--1980}. Cambridge University Press.
\end{thebibliography}

\end{document}
